How old can a predicted action be before it should be refreshed? Action-chunk
policies forecast a trajectory from each observation, yet execution rules
apply one temporal setting to the action vector. We diagnose component
dependence by taking arm and gripper commands from ACT
chunks while aligning both to the same physical control step.
These probes query the policy once per step and are diagnostic
interventions, not deployment schedules. In a preregistered 140-block ACT
evaluation, moving a one-second-old prediction from the gripper to the arm
changed success from 83/140 to 38/140, a 32.14-point difference. On a separate
development cohort, the fresh-arm/stale-gripper branch remained above the
stale-arm/fresh-gripper branch from 0.10 to 1.60~s, with observed separations of
21.43--54.76 points. Translation staleness was also substantially more damaging
than rotation staleness. In executable schedules, extending gripper commitment
improved success by 4.667 and 5.778 points at two arm cadences with nearly
matched policy-query rates. Coherent H16 nevertheless remained the strongest
operating point. Same-target disagreement ranked the arm
components in the opposite order from behavior, and the probe couples source
age with prediction lookahead. The mechanism and generality beyond the
evaluated ACT policies therefore remain unresolved.
